% Pursuing Stacks (cotes 134-2 to 134-8): archive cote 134-6, site page 59 (PDF page 60), typescript page 377.
% typescript. Draft compared: ps5.tex:4325-4397.
% AI-transcribed by Claude Opus 5.5 (claude-opus-5-5, adaptive thinking, effort high), Message Batches API, batch ps3 (msgbatch_01XdFmuVVBAcgc7EFxdtCqsj), 2026-11-12; not yet checked by eye.
\scanpage{134-6/59}
\mspage{377}

expresses cointegration $Rf_*$ with coefficients \xout{in coefficients on B' comin}
coming from downstairs, namely $f^*(K^\bullet)$, where $K^\bullet$ is any differential \xout{compl}
complex in $B^\wedge_{ab}$ with degrees bounded from below (thus defining an object
in $D^+(B^\wedge_{ab})$). The relevant formula is
\[ (12) \qquad Rf_*(\xout{F}f^*(K^\bullet)) \simeq \underline{Hom}^{\bullet\bullet}(L_\bullet^{B'/B} , K^\bullet) \quad , \]
an isomorphism in $D^+(B^\wedge_{ab})$, where $\underline{Hom}^{\bullet\bullet}$ designates the double complex
in $B^\wedge_{ab}$ obtained by applying $\underline{Hom}$ componentwise, more accurately the
associated simple complex, and where $\underline{Hom}$ is the internal $\underline{Hom}$ in the
category $B^\wedge_{ab}$ namely the \xout{sheafxx} (pre)sheaf of additive homomorphisms of
a given abelian (pre)sheaf ($L_n$ say) into another ($K^m$ say). The proof
of (12) is essentially trivial, it is just the computational interpreta-%
tion, in terms of using projective resolutions, of the adjunction formu-%
la ``localized on B''
\margin{\textcircled{$Lf^{ab}_!$}}
\[ \xout{(13)}\unsure{(12')} \qquad Rf_*(Lf^*(K^\bullet)) \hand{\simeq} R\underline{Hom}(Lf^{ab}_!(\underline{Z}_{B'}), K^\bullet) \quad , \]
which is a particular case of the more general \inserted[typed]{``adjunction} formula''
\margin{\textcircled{$F'_\bullet$}}
\[ (13) \qquad Rf_*(R\underline{Hom}(F'_\bullet , Lf^*(K^\bullet)) \hand{\simeq} R\underline{Hom}(Lf^{ab}_!(F'_\bullet), K^\bullet) \quad , \]
valid for
\[ F'_\bullet \text{ in } D^-(B'^\wedge_{ab}) \ , \quad K^\bullet \text{ in } D^+(B^\wedge_{ab}) \quad , \]
(12') following from (13) by taking $F'_\bullet = \underline{Z}_{B'}$ .

\uline{Remarks.}
We may view (13), and it's particular case (12) or (12'), as
the main formula relating the \uline{homology} and \uline{cohomology} invariants \xout{ax}\inserted[typed]{for} a map
f in (Cat), or equivalently, the \inserted[typed]{(abelian)} \uline{integration} and \uline{cointegration} operati-%
ons defined by f. It now occurs to me that this formula, \xout{andx} and the
variance formalism in which it inserts, is valid more generally whenever
we have a map f between two \xout{topox} ringed topoi, such that $f_!$ exists
for sheaves of sets, hence there exists too a corresponding functor
$f^{mod}_!$ for sheaves of modules. The fact that we have been restricting to
the case of the constant sheaves of rings defined by $\underline{Z}$ is'nt relevant,
and (in the case of topoi defined by objects in (Cat), \deleted{\unsure{the}}\inserted{hence} with suffi-%
ciently many projective sheaves of sets) the formalism of the subcate-%
gories Add(B) and Addinf(B) in $B^\wedge_{ab}$ can be generalized equally to arbit-%
rary \xout{commutative} sheaves of rings on B. At present I do'nt see though
any \xout{interesting}\inserted[typed]{striking} particular case where this generalization would seem
useful.

\margin{59}
